% ===== PRÉAMBULE CANONIQUE — à coller VERBATIM en tête de CHAQUE .tex =====
\documentclass[11pt,a4paper]{article}
\usepackage[utf8]{inputenc}\usepackage[T1]{fontenc}\usepackage{lmodern}
\usepackage[french]{babel}
\usepackage{amsmath,amssymb,mathtools}\usepackage{icomma}
\usepackage{siunitx}\sisetup{locale=FR}  % \qty{}{}, \si{}, \ang{}
\usepackage{xcolor}\usepackage{colortbl}
\usepackage{tikz}\usepackage{pgfplots}\pgfplotsset{compat=1.18}
\usepackage[siunitx,europeanresistors]{circuitikz} % circuits (élec.)
\usepackage[most]{tcolorbox}\usepackage{enumitem}\usepackage{array,booktabs}
\usepackage[a4paper,margin=2.2cm]{geometry}
\usetikzlibrary{arrows.meta,calc,angles,quotes,positioning,patterns,%
  backgrounds,shapes.geometric,decorations.pathmorphing,decorations.markings,intersections}
\definecolor{primary}{RGB}{222,49,99}\definecolor{primarydark}{RGB}{150,22,55}
\definecolor{defcol}{RGB}{0,114,178}\definecolor{methcol}{RGB}{52,92,148}
\definecolor{excol}{RGB}{0,135,90}\definecolor{attncol}{RGB}{200,40,55}
\tcbset{boxbase/.style={enhanced,breakable,boxrule=0.9pt,arc=2.5pt,
  left=3mm,right=3mm,top=2.3mm,bottom=2.3mm,fonttitle=\bfseries,coltitle=white}}
% --- boîtes TUTORIEL (noms EXACTS, ne pas renommer) ---
\newtcolorbox{cle}[1][]{boxbase,colback=defcol!6,colframe=defcol,
  title={\textsc{À retenir}\ifx\relax#1\relax\relax\else\textmd{ : \itshape #1}\fi}}
\newtcolorbox{methode}[1][]{boxbase,colback=methcol!7,colframe=methcol,sharp corners,
  title={\textsc{Méthode}\ifx\relax#1\relax\relax\else\textmd{ --- \itshape #1}\fi}}
\newtcolorbox{exemplebox}[1][]{boxbase,colback=excol!5,colframe=excol,
  title={\textsc{Exemple}\ifx\relax#1\relax\relax\else\textmd{ : \itshape #1}\fi}}
\newtcolorbox{piege}[1][]{boxbase,colback=attncol!6,colframe=attncol,
  title={\textsc{Piège}\ifx\relax#1\relax\relax\else\textmd{ : \itshape #1}\fi}}
\newtcolorbox{remarque}[1][]{boxbase,colback=black!4,colframe=black!45,
  title={\textsc{Remarque}\ifx\relax#1\relax\relax\else\textmd{ : \itshape #1}\fi}}
% --- boîtes EXERCICE / CORRIGÉ (noms EXACTS, auto-comptées) ---
\newtcolorbox[auto counter]{exo}[1][]{boxbase,colback=primary!4,colframe=primary,
  title={\textsc{Exercice}~\thetcbcounter\ifx\relax#1\relax\relax\else\textmd{ --- \itshape #1}\fi}}
\newtcolorbox[auto counter]{corrige}[1][]{boxbase,colback=excol!4,colframe=excol,
  title={\textsc{Corrigé}~\thetcbcounter\ifx\relax#1\relax\relax\else\textmd{ --- \itshape #1}\fi}}
\newcommand{\vv}[1]{\vec{#1}}\newcommand{\ds}{\displaystyle}
\newcommand{\R}{\mathbb{R}}\newcommand{\N}{\mathbb{N}}\newcommand{\Z}{\mathbb{Z}}
% =========================================================================
\begin{document}
\sisetup{per-mode=symbol}
% ================= MOYEN =================

\begin{exo}[vitesse et période à \qty{1600}{\kilo\metre}]
Un satellite orbite à $r=R_T+h=\qty{8.0e6}{\metre}$ ($\mathcal{G}M_T=\qty{4.0e14}{}$ SI).
\begin{enumerate}[label=\alph*)]
  \item Calcule sa vitesse $v=\sqrt{\mathcal{G}M_T/r}$.
  \item Calcule sa période $T=\dfrac{2\pi r}{v}$ (en secondes, puis en minutes).
\end{enumerate}
\end{exo}

\begin{exo}[la période par la loi de Kepler]
Un satellite gravite à $r=\qty{1.0e7}{\metre}$. En utilisant
$T^2=\dfrac{4\pi^2 r^3}{\mathcal{G}M_T}$ (avec $4\pi^2\approx 39{,}5$ et
$\mathcal{G}M_T=\qty{4.0e14}{}$ SI), calcule la période $T$.
\end{exo}

\begin{exo}[le satellite géostationnaire]
Un satellite géostationnaire a une période $T=\qty{86400}{\second}$ et gravite à
$r=\qty{4.22e7}{\metre}$ du centre.
\begin{enumerate}[label=\alph*)]
  \item Calcule sa vitesse angulaire $\omega=\dfrac{2\pi}{T}$.
  \item Calcule sa vitesse linéaire $v=\omega\,r$.
\end{enumerate}
\end{exo}

\begin{exo}[comparer deux vitesses]
Le satellite $A$ est à $r_A=\qty{8.0e6}{\metre}$, le satellite $B$ à
$r_B=\qty{2.0e7}{\metre}$.
\begin{enumerate}[label=\alph*)]
  \item Établis l'expression du rapport $\dfrac{v_A}{v_B}$ en fonction de $r_A$ et $r_B$.
  \item Calcule sa valeur. Lequel va le plus vite ?
\end{enumerate}
\end{exo}

\begin{exo}[deux périodes par Kepler]
Le satellite $2$ orbite à un rayon $8$ fois plus grand que le satellite $1$
($r_2=8\,r_1$).
\begin{enumerate}[label=\alph*)]
  \item À partir de $T^2\propto r^3$, établis l'expression du rapport $\dfrac{T_2}{T_1}$.
  \item Calcule sa valeur numérique.
\end{enumerate}
\end{exo}

\end{document}
